\documentclass[11pt]{article}
\usepackage[a4paper,margin=25mm]{geometry}
\usepackage{amsmath,amssymb,amsthm}
\usepackage[hidelinks]{hyperref}
\newtheorem{theorem}{Theorem}
\begin{document}
\title{Finite samples have zero-dimensional Zariski closure}
\author{Chengzhe Gao}\date{}\maketitle
Conjecture 2839 claims that the Zariski closure of a generic $n$-point
sample in ambient dimension $d$ has dimension $\min(n-1,d)$.
That formula fails already for $n=2,d=1$, for \emph{every} two-point
sample, so choosing a generic sample cannot repair it.

\begin{theorem}
For any two distinct points $a,b$ on the affine line over an algebraically
closed field, their Zariski closure has dimension zero.
\end{theorem}
\begin{proof}
The polynomial $(X-a)(X-b)$ has zero set exactly $\{a,b\}$, because the
field has no zero divisors. Thus the sample itself is Zariski closed and
is its own Zariski closure. Each singleton is also closed, being the zero
set of $X-a$ or $X-b$.

A nonempty irreducible closed subset of $\{a,b\}$ must be contained in
one of those singletons: apply irreducibility to the cover
$\{a,b\}=\{a\}\cup\{b\}$. Hence the only nonempty irreducible closed
subsets are the two singletons. There is a chain of length zero and no
strict chain of length one. By the standard dimension definition using
chains of nonempty irreducible closed subsets, the dimension is exactly
zero. The asserted formula instead gives $\min(2-1,1)=1$.
\end{proof}

The proof is uniform in $a,b$, including points in every nonempty generic
locus of the configuration space. It does not substitute a specially
chosen rational sample for a generic one. The affine span of two distinct
points in the affine line does have dimension one, but affine span is a
different operation from the Zariski closure named in both language
versions of the source. No finite restriction on polynomial degree is
part of the ordinary Zariski closure.

\paragraph{Formal verification.}
\texttt{DomainData} lists the ordinary integral-domain laws, all valid
over any field; the main theorem is universally quantified in those laws.
An integer model separately checks consistency without restricting that
universal theorem. \texttt{Polynomial} contains precisely constant,
variable, addition, multiplication and negation expressions. Every usual
coefficient-list polynomial is embedded by Horner's construction, with
a proved evaluation bridge.

Zariski closed sets are zero sets of arbitrary collections of these
polynomials. Closure is the actual common-zero set $V(I(S))$ of all
polynomials vanishing on the sample. The file proves inclusion, closedness,
and the least-closed-superset property, then identifies the two-point
closure exactly. Irreducibility quantifies over arbitrary closed covers.
Dimension at least $k$ is witnessed by a length-$k$ strict chain of
nonempty irreducible closed subsets. The proof classifies all such subsets
inside the two-point set and excludes every chain of length one, while
constructing one of length zero.

The final theorem rules out \emph{any} distinct pair with the dimension
predicted at $n=2,d=1$. It therefore excludes any putative generic
instance as well. Lean 4.19.0 verifies the complete polynomial, closure,
irreducibility, and dimension chain with only \texttt{Std}, no custom
axioms, no admitted declarations, and no native evaluation proof.

\end{document}
